% =====================================================================
%  A benchmark cannot order more methods than its labels can distinguish
%  Manuscript for Nature Methods (Article)
%
%  Compiles with pdfLaTeX on a stock TeX Live / Overleaf installation.
%  No journal class file is required: Nature Portfolio accepts any
%  reasonable format at initial submission and asks only for double
%  spacing, continuous line numbers and page numbers.
%
%      pdflatex main && bibtex main && pdflatex main && pdflatex main
% =====================================================================
\documentclass[11pt,a4paper]{article}

\usepackage[T1]{fontenc}
\usepackage[utf8]{inputenc}
\usepackage{lmodern}
\usepackage[margin=1in]{geometry}
\usepackage{amsmath,amssymb}
\usepackage{graphicx}
\usepackage{booktabs}
\usepackage{array}
\usepackage{caption}
\usepackage{setspace}
\usepackage{lineno}
\usepackage{xcolor}
\usepackage[hidelinks]{hyperref}
\usepackage{authblk}
\usepackage{textcomp}

% ---- Unicode the sources use, declared so nothing silently fails --------
\DeclareUnicodeCharacter{03B1}{\ensuremath{\alpha}}
\DeclareUnicodeCharacter{03B2}{\ensuremath{\beta}}
\DeclareUnicodeCharacter{03B3}{\ensuremath{\gamma}}
\DeclareUnicodeCharacter{03B4}{\ensuremath{\delta}}
\DeclareUnicodeCharacter{03B5}{\ensuremath{\varepsilon}}
\DeclareUnicodeCharacter{03BA}{\ensuremath{\kappa}}
\DeclareUnicodeCharacter{03BB}{\ensuremath{\lambda}}
\DeclareUnicodeCharacter{03BC}{\ensuremath{\mu}}
\DeclareUnicodeCharacter{03BD}{\ensuremath{\nu}}
\DeclareUnicodeCharacter{03C1}{\ensuremath{\rho}}
\DeclareUnicodeCharacter{03C3}{\ensuremath{\sigma}}
\DeclareUnicodeCharacter{0394}{\ensuremath{\Delta}}
\DeclareUnicodeCharacter{03A3}{\ensuremath{\Sigma}}
\DeclareUnicodeCharacter{2264}{\ensuremath{\leq}}
\DeclareUnicodeCharacter{2265}{\ensuremath{\geq}}
\DeclareUnicodeCharacter{2260}{\ensuremath{\neq}}
\DeclareUnicodeCharacter{2248}{\ensuremath{\approx}}
\DeclareUnicodeCharacter{00D7}{\ensuremath{\times}}
\DeclareUnicodeCharacter{00B7}{\ensuremath{\cdot}}
\DeclareUnicodeCharacter{00B2}{\ensuremath{^2}}
\DeclareUnicodeCharacter{2192}{\ensuremath{\rightarrow}}
\DeclareUnicodeCharacter{21D2}{\ensuremath{\Rightarrow}}
\DeclareUnicodeCharacter{2308}{\ensuremath{\lceil}}
\DeclareUnicodeCharacter{2309}{\ensuremath{\rceil}}
\DeclareUnicodeCharacter{230A}{\ensuremath{\lfloor}}
\DeclareUnicodeCharacter{230B}{\ensuremath{\rfloor}}
\DeclareUnicodeCharacter{2208}{\ensuremath{\in}}
\DeclareUnicodeCharacter{2286}{\ensuremath{\subseteq}}
\DeclareUnicodeCharacter{2229}{\ensuremath{\cap}}
\DeclareUnicodeCharacter{222A}{\ensuremath{\cup}}
\DeclareUnicodeCharacter{0394}{\ensuremath{\Delta}}
\DeclareUnicodeCharacter{1D53C}{\ensuremath{\mathbb{E}}}
\DeclareUnicodeCharacter{2026}{\dots}
\DeclareUnicodeCharacter{2014}{---}
\DeclareUnicodeCharacter{2013}{--}
\DeclareUnicodeCharacter{2011}{-}
\DeclareUnicodeCharacter{2019}{'}
\DeclareUnicodeCharacter{201C}{``}
\DeclareUnicodeCharacter{201D}{''}
\DeclareUnicodeCharacter{2265}{\ensuremath{\geq}}
\DeclareUnicodeCharacter{2205}{\ensuremath{\emptyset}}
\DeclareUnicodeCharacter{221E}{\ensuremath{\infty}}
\DeclareUnicodeCharacter{2211}{\ensuremath{\sum}}
\DeclareUnicodeCharacter{03C0}{\ensuremath{\pi}}
\DeclareUnicodeCharacter{00B0}{\ensuremath{^\circ}}
\DeclareUnicodeCharacter{03A9}{\ensuremath{\Omega}}
\DeclareUnicodeCharacter{03B9}{\ensuremath{\iota}}
\DeclareUnicodeCharacter{03C9}{\ensuremath{\omega}}
\DeclareUnicodeCharacter{1D9C}{\ensuremath{^{c}}}
\DeclareUnicodeCharacter{2080}{\ensuremath{_0}}
\DeclareUnicodeCharacter{2081}{\ensuremath{_1}}
\DeclareUnicodeCharacter{211D}{\ensuremath{\mathbb{R}}}
\DeclareUnicodeCharacter{2209}{\ensuremath{\notin}}
\DeclareUnicodeCharacter{2212}{\ensuremath{-}}
\DeclareUnicodeCharacter{2218}{\ensuremath{\circ}}
\DeclareUnicodeCharacter{2227}{\ensuremath{\wedge}}

\captionsetup{labelfont=bf,font=small,justification=justified,
              singlelinecheck=false}
\renewcommand{\arraystretch}{1.15}
\newcommand{\eps}{\varepsilon}
\newcommand{\code}[1]{\texttt{\small #1}}

\sloppy
\emergencystretch=3em
\doublespacing
\linenumbers

\title{\bfseries A benchmark cannot order more methods than its labels
can distinguish}

\author[1]{Tommaso R. Marena}
\affil[1]{Correspondence: \texttt{marenatommaso@gmail.com}}
\date{}

\begin{document}
\maketitle

% ---------------------------------------------------------------------
\begin{abstract}
\noindent
Community benchmarks rank methods on imperfect labels, and the resulting order
is treated as data. We prove that a benchmark with annotation error rate
$\eps$ can place at most $\lceil 1/2\eps \rceil$ methods in a certifiable
order --- however many enter, and however analysed --- and that beyond this
the ordering is not a function of the observations. Measured from
repeat determinations of the same protein, $\eps$ gives the Critical Assessment
of Intrinsic Disorder (CAID) a capacity of seven; it admits 117. The resolution
is discarded rather than absent: a ranking statistic is corrupted only when two
labels flip in opposite directions, an exactly second-order event, so scoring
within-protein pairs raises the capacity to 51. We specify that protocol,
re-score CAID3 under it, and show the ranking reproduces on a held-out round as
well as the published one. The same argument governs how a screen should be
reported. All theorems are machine-checked in Lean~4.
\end{abstract}

\vspace{1em}

% =====================================================================
\section*{Introduction}
% =====================================================================

Intrinsically disordered regions are not folded, and the experiment that
reports them --- a deposited structure with residues absent from the
coordinates --- reports a property of one crystal, one construct, one
condition. The Critical Assessment of Intrinsic Disorder (CAID)\cite{caid1,
caid2} turns those annotations into a leaderboard: every residue of every
target is pooled into one area under the ROC curve, and 117 submissions are
placed in order.

\textbf{That number is mostly not the question it is read as.} A pooled AUC is the
probability that a random disordered residue outranks a random ordered one,
and every such pair is either within one protein or between two. The statistic
therefore decomposes exactly,
\[
  \mathrm{AUC}_{\text{pooled}}
    = w_{\text{within}}\,\mathrm{AUC}_{\text{within}}
    + w_{\text{between}}\,\mathrm{AUC}_{\text{between}},
\]
and on CAID3 the within-protein weight is 0.51\%--2.9\%. That much is
arithmetic --- for $n$ proteins of comparable size about $1/n$ of pairs fall
inside one --- and licenses only an empirical question: does the metric's 1:200
weighting change who wins? On three of five CAID3 benchmarks it does.

\textbf{And it is being asked to resolve more than it can.} Every benchmark
orders methods using annotations, and annotations are wrong at some rate. How
many methods that rate permits to be ordered \emph{at all} is a question the
field does not ask. We prove the answer is $\lceil 1/2\eps \rceil$, prove it is
attained, and prove the converse: past that count, two close methods admit two
ground truths consistent with every observation, one favouring each. The only
remedies are better labels --- or a different statistic. This paper supplies the
second: the bound constrains \emph{labels}, AUC is a function of ordered
\emph{pairs}, and the two rates are not the same quantity.

% =====================================================================
\section*{Results}
% =====================================================================

\subsection*{CAID's statistic is 97--99.5\% a between-protein question}

The decomposition above is exact, and its first term has a property that makes
it more than bookkeeping. \code{auc\_within\_strictMono\_invariant} (Lean~4,
\code{sorry}-free): $\mathrm{AUC}_{\text{within}}$ is invariant under any
per-protein strictly monotone recalibration. Rescore every chain by its own
exponential, cube or logistic squash and the quantity does not move.
\code{auc\_pooled\_shift\_diff} completes it: the entire effect of such a
recalibration on the reported number falls in the between-protein term. The
split is therefore a partition of CAID's statistic into the part no
chain-by-chain rescoring can touch and the part that is nothing else
(Fig.~\ref{fig:decomp}a; Supplementary Table~S1).

On three of five CAID3 benchmarks the method the challenge declares best is
21st to 25th at telling which residues of a chain are disordered
(Fig.~\ref{fig:decomp}b): fl\-DPnn3a is 25th of 57 on Disorder-NOX,
UdonPred-combined 21st of 69 on Binding, LINKER-Pred2 25th of 90 on Linker. The
three Linker specialists sweep the pooled table at \#1, \#2 and \#3 and sit at
\#25, \#26 and \#29 within protein: excellent at ranking whole chains by linker
content, unremarkable at locating linkers inside one.

Each disagreement is a certificate rather than an observation.
\code{inversion\_requires\_between\_gap} states that the pooled order reverses
the within-protein order only when $w_{\text{within}}\times$(within gap)
$< w_{\text{between}}\times$(between gap), so an inversion \emph{forces} a
between-protein gap of at least $(w_{\text{within}}/w_{\text{between}})$ times
the within-protein gap. On Disorder-NOX every certificate holds and each is
over-determined by 562--2305$\times$ (Fig.~\ref{fig:decomp}c): a within-protein
deficit of 0.0218 requires a between-protein advantage of 0.000148 to be
overturned, and the one present is 562 times larger. Calibration differences of
ordinary size do not merely outweigh discrimination differences of ordinary
size; they swamp them by three orders of magnitude.

The counterweight belongs here and not in a footnote: Spearman(pooled, within)
is $+0.964$ on Disorder-PDB and never below $+0.849$, so the decomposition does
not overturn CAID3 (Fig.~\ref{fig:rescored}a). What it does is locate where the
two orderings disagree, and the movement is large where they do.

\subsection*{A benchmark can order at most $\lceil 1/2\eps\rceil$ methods}

\code{BenchmarkCapacity.card\_le\_benchCapacity} (Lean~4, \code{sorry}-free;
axioms \code{propext}, \code{Classical.choice}, \code{Quot.sound}). A benchmark
with $n$ targets, scores averaged into $[0,1]$ and therefore on the grid of
denominator $n$, can place in a certified order at most
\[
  k(n,\eps,\delta) \;=\; \frac{n}{\lfloor c\,n\rfloor + 1} + 1,
  \qquad c = \max(\delta,\, 2\eps)
\]
methods, \emph{however many are entered}. The factor two is the two-sided price
of label noise; $\delta$ is the smallest difference worth calling a difference
($\delta = 0$ throughout, and any positive $\delta$ only lowers $k$).
\code{benchCapacity\_attained} exhibits a family of exactly $k$ methods the
benchmark does resolve, in each of three score models, so $k$ is the capacity
and not an estimate. \code{benchCapacity\_noise\_only} gives the $n$-free
ceiling $k \le \lceil 1/2\eps \rceil$: collecting more targets does not buy
resolution the labels lack.

$\eps$ is measured, not assumed. MobiDB\cite{mobidb} publishes the
missing-residue call of each deposited structure, up to 806 for one protein.
Across 160 CAID3 Disorder-PDB targets, 4{,}886 of 61{,}013 evidenced residues
are context-dependent --- missing in some structures and observed in others ---
giving $\eps = 0.0801$; an independent estimator over all pairs of structures of
the same protein gives 0.0651. That is not measurement error but the fact that
the label is not a function of the sequence: a region ordered in one crystal
form and not another is scored as one number per residue against a reference
that depends on which structure was consulted.

Three independent routes --- the target grid, the residue-level form
\code{card\_le\_capacity\_labelNoise}, and the $n$-free ceiling --- give a
capacity of \textbf{seven} on every benchmark of two CAID rounds
(Fig.~\ref{fig:capacity}a; Supplementary Table~S6). CAID3 admits 117 entrants
and can place at most seven of them in a certified order. To order the full
field, the annotation error rate would have to fall to 0.427\%: a nineteen-fold
improvement in label quality, not a larger benchmark.

\textbf{This is impossibility, not low power.} \code{unresolvable\_pair} is the
converse and it is what separates this from a power calculation. A benchmark
observes the annotation, not the truth, and every labelling within the noise
budget is a consistent truth. For two methods whose measured scores are close,
the proof \emph{constructs two such truths}, one making each strictly better,
both compatible with everything the benchmark recorded. The ordering is
therefore not a function of the data, and no analysis of those data recovers
it: not a larger bootstrap, not permutation testing, not leave-one-out, not
Monte Carlo. \code{over\_capacity\_has\_close\_pair} closes the loop: once the
entry list exceeds capacity, such a pair necessarily exists.

Our own results are consistent with this, which is the point rather than an
excuse. Five pre-registered architectures and an ensemble each failed to
separate from PUNCH2 on pooled AUC, and each of the five also failed the
endpoint its own registration named against a regime-matched control, with none
of fifteen paired comparisons surviving Holm correction (Supplementary
Note~S4). We were trying to resolve a pair the benchmark cannot resolve.

\subsection*{Pairwise scoring makes the noise second-order}

The bound constrains \emph{labels}. AUC is a function of ordered \emph{pairs},
and the noise a ranking statistic suffers is the rate at which the annotation
\emph{reverses} a pair.

Let $T$ be the truth on one protein and $L$ the annotation. A pair $(p,q)$ is
discordant when $T$ orders it one way and $L$ the other, which requires
$T(p){=}1, L(p){=}0$ \emph{and} $T(q){=}0, L(q){=}1$: both residues flip, in
opposite directions. With $d = |T \setminus L|$ and $u = |L \setminus T|$,
\[
  \mathrm{discordant}(T,L) = 2\,d\,u \quad\text{exactly},
  \qquad \mathrm{noise}(T,L) = d+u ,
\]
and by AM--GM $\mathrm{discordant} \le \mathrm{noise}^2/2$. Both hold with no
hypotheses (\code{discordant\_eq\_flip\_product},
\code{discordant\_le\_noise\_sq}). The pairs a pairwise protocol can score are
those \emph{both} labellings order, counted exactly as well:
$|\mathrm{comparable}(T,L)| = 2(|T \cap L|\cdot|(T \cup L)^{c}| + du)$
(\code{card\_comparablePairs}), so the pairwise noise rate is exactly
$\eps_{\text{pair}} = du/(|T \cap L|\cdot|(T \cup L)^{c}| + du)$.
\code{noise\_orderKey} closes the loop: on comparable pairs that quantity
\emph{is} the label noise of the answer key the annotation induces, so the
capacity theorem applies at the pair level verbatim
(\code{pairwise\_capacity\_bound}).

Measured on 2{,}746 pairs of deposited structures of the same protein, across
the 159 CAID3 Disorder-PDB targets with at least two (Fig.~\ref{fig:capacity}b;
Supplementary Table~S7):

\begin{center}\small
\begin{tabular}{lrr}
\toprule
 & pooled & median protein \\
\midrule
label-flip rate $\eps_{\text{label}}$ & 0.0651 & 0.0342 \\
pairwise discordance $\eps_{\text{pair}}$ & \textbf{0.0100} & 0.0006 \\
capacity $\lceil 1/2\eps \rceil$ & \textbf{8 $\rightarrow$ 51} & \\
\bottomrule
\end{tabular}
\end{center}

\noindent
A benchmark scored on within-protein pairwise comparisons can order 51 methods
where the residue-level protocol orders 8. CAID3 has 117 entrants.

\textbf{The closed form does not apply here, and the theorem is what says so.}
\code{nuPair\_le\_two\_eps\_sq} gives $\eps_{\text{pair}} \le 2\eps^2$, and with
it a capacity quadratic in label quality rather than linear --- but it requires
the \emph{agreement} classes to balance, $|T \cap L| = |(T \cup L)^{c}|$. That
hypothesis is load-bearing: the denominator collapses when the classes are
lopsided, and the gain collapses with it. On a reference that is 31.6\%
disordered it fails comprehensively --- pooled, the agreement classes stand at
240{,}506 disordered against 555{,}402 ordered --- and it holds on 32 of 2{,}746
structure pairs (Fig.~\ref{fig:capacity}c). On all 32 of those the bound holds;
pooled, the measured rate exceeds it, 0.00996 against 0.00848. The sixfold gain
is therefore a \emph{measurement}, not a prediction from the label-noise rate,
and the protocol below says \emph{estimate} $\eps_{\text{pair}}$ rather than
derive it. The two results the argument rests on --- the identity and the exact
denominator --- carry no hypotheses; the closed form that would let a benchmark
skip the measurement is exactly the part CAID3 cannot use.

\subsection*{A protocol, and the field re-scored}

\textbf{The CAID pairwise protocol} (Methods) scores a method by the
\emph{unweighted} mean of its per-target Mann--Whitney AUCs, over the targets
carrying both classes, among methods that predicted every evaluated residue of
every target. Two methods are reported as ordered only when a paired Wilcoxon
test over targets rejects equality after Holm correction, and the number of
methods the reference can order at all is
$\lceil 1/2\eps_{\text{pair}} \rceil$, with anything beyond that reported as an
unresolved group rather than as a rank.

The unweighted mean is what makes the protocol calibration-invariant --- each
per-target AUC is unchanged by any strictly monotone recalibration of that
target's scores, so any function of them is too --- and the capacity step is what
keeps it honest: a rank the labels cannot support is not printed.

Re-scoring CAID3 (Fig.~\ref{fig:rescored}; Supplementary Table~S3), the two
orderings agree overall ($\rho = 0.844$ over 357 method--benchmark pairs) and
disagree sharply in places. On Disorder-PDB PUNCH2 falls from 4th to 7th; on
Disorder-NOX \textbf{AlphaFold-pLDDT moves from 39th to 6th} --- a training-free
confidence score, never trained on disorder, sixth in the field at the
residue-level question and thirty-ninth on the number the challenge reports.

The protocol also declares what it cannot order. The leader is separated at the
Holm-corrected margin from 55 of 59 eligible methods on Disorder-PDB, 51 of 59
on Disorder-NOX, 56 of 71 on Binding, 50 of 71 on Binding-IDR and 87 of 92 on
Linker (Fig.~\ref{fig:rescored}c), against seven orderings in total for the
residue-level protocol. Those counts exceed the capacity of 51 and are allowed
to: the capacity bounds a certified \emph{totally ordered family}, whereas 87
separations against one leader are 87 statements about the pairs
$\{\text{leader}, X\}$, which need not compose into a chain of 88. Step 6
governs the ranking the protocol prints, not the number of pairs a test rejects.

\textbf{The extra resolution is not noise.} Fifty-seven methods entered both
CAID2 and CAID3; ranking them on CAID3 under each protocol and asking which
reproduces their CAID2 ranking gives $\rho = 0.997$ for pairwise against 0.976
for pooled on Disorder-PDB, and 0.83--0.997 against 0.82--0.95 across four
benchmarks (Fig.~\ref{fig:rescored}d). The pairwise ranking transfers as well as
the pooled one. It is \emph{not} shown to be more valid and we do not claim it
is: three of four benchmarks show no significant difference and one favours the
pooled protocol. The case rests on capacity, not on a validity advantage the
data do not show.

CAID4's references are not yet published, so this protocol can be adopted rather
than applied retrospectively, which is the difference between a critique and a
contribution.

\subsection*{An operating guarantee, and the price of it}

Every disorder predictor in CAID reports a ranking statistic. None answers the
question a biologist has: \emph{can I act on this call, and how often will it be
wrong?} Conformal risk control\cite{crc} supplies the missing guarantee,
assuming nothing about the model, the distribution, or the calibration of the
scores: for a bounded loss monotone in a threshold it returns a $\lambda$ with
$\mathbb{E}[L_{\text{test}}(\lambda)] \le \alpha$ over a fresh \emph{protein},
the unit the split respects (Methods).

Applied to all 70 full-coverage CAID3 Disorder-PDB methods
(Fig.~\ref{fig:operating}a): validity is free and identical for everyone --- the
threshold is chosen to make it so --- and the entire content is the price, which
runs from 32.0\% of the protein flagged (DisorderNet) to 97.3\%. A second rule
buys the identical guarantee by flagging a fixed \emph{quantile of each
protein's own scores}, which depends only on the ordering inside a chain and is
therefore invariant under per-protein recalibration. The two costs price the two
abilities the decomposition separates --- a global threshold prices
discrimination \emph{and} calibration, a per-protein quantile prices
discrimination alone --- and their difference is what protein-level calibration
is worth operationally (Fig.~\ref{fig:operating}b). The leaders lean on it twice
as hard as we do, which yields the sharpest comparison in the paper
(Fig.~\ref{fig:operating}c): on Disorder-PDB, pooled AUC separates DisorderNet
from PUNCH2 by $+0.0043$ at $p = 0.20$ --- a difference we have never been able
to call real. The calibration-invariant operating cost separates them by 5.3
points (55.1\% against 60.4\% of the protein flagged), and by 10.2 points on
Disorder-NOX. \emph{The two methods are statistically inseparable on the
benchmark's own metric and are not close as instruments.}

\subsection*{Reporting a screen: the unit of testing is a design choice too}

The capacity result says the unit of \emph{scoring} is a choice. The same
structure appears one step downstream, when a predictor is run over a proteome.

A screen's candidates are not independent: they share a calibration run, a model
and a training set. Under arbitrary dependence the Benjamini--Hochberg
procedure\cite{bh} needs the harmonic correction\cite{by}. We prove it from
scratch in a finite, measure-theory-free setting, assuming only that the p-value
of a true null is valid: \code{selfConsistent\_fdr\_le\_harmonic} bounds
$\mathbb{E}[\mathrm{FDP}]$ by $q H_m |H_0|/m$ for \emph{any} self-consistent
step-up rule under an \emph{arbitrary} joint law. The mechanism, isolated as
\code{wgt\_decomp}, is why the joint law never enters: the weight $1/|R|$ of a
discovery telescopes into \emph{nested} events $\{i \in R, |R| \le j\}$, each
of which self-consistency contains in the single-candidate event
$\{p_i \le jq/m\}$, which validity alone bounds. The price is logarithmic,
$\log(m{+}1) \le H_m \le 1 + \log m$, against Bonferroni's $m$
(Fig.~\ref{fig:screen}a).

\textbf{The correction is necessary, not conservative.} For every screen size
and level we construct an explicit joint law in which all $m$ hypotheses are
null, every p-value is valid, BH rejects exactly a cyclic window, and therefore
$\mathbb{E}[\mathrm{FDP}] = qH_m$ \emph{exactly}
(\code{bh\_fdr\_eq\_harmonic}). The factor is attained and cannot be lowered;
uncorrected BH strictly exceeds its nominal level from two candidates on
(Fig.~\ref{fig:screen}c).

The factor therefore depends on one thing only: how many hypotheses the screen
states. And disorder is a property of a region. \code{fdp\_lift} shows that a
region-level report and its residue-level reading have the \emph{same} false
discovery proportion --- both numerator and denominator scale by the region
length --- so \code{region\_screen\_fdr\_le} gives residue-level control at
$\alpha$ from a region-level procedure at $\alpha/H_M$, and
\code{harmonic\_gain\_log} bounds the saving below by $\log b - 1$. On the CAID
references a region-level screen may test at a threshold 1.61--2.05$\times$
larger for the same residue-level rate (Fig.~\ref{fig:screen}b).

We ran the screen. Conformal p-values calibrated on the ordered blocks of 80\%
of chains, Benjamini--Yekutieli at $\alpha/H_M$ on the rest, 50 splits: at
$\alpha = 0.10$ with 100-residue blocks the region-level screen reports 2{,}942
residues at a realised false discovery proportion of 0.024 and power 0.568,
against 203 residues and power 0.036 for the residue-level screen at the same
guaranteed rate --- \textbf{14.5$\times$ the discoveries}. A third constraint,
not visible in the theory, dominates both: a conformal p-value cannot fall below
$1/(n_{\text{cal}}{+}1)$, and a step-up rule at level $\alpha/H_m$ over $m$
hypotheses applies thresholds as small as $\alpha/(mH_m)$, so a screen is
arithmetically incapable of a discovery unless
$n_{\text{cal}} + 1 \ge m H_m / \alpha$. That requirement is \emph{linear} in
the number of hypotheses. Stating one per region rather than one per residue
reduces it by 13--178$\times$ on these references, which is a far larger effect
than the harmonic factor itself, and it is the constraint that binds first.

\subsection*{DisorderNet}

DisorderNet is a frozen-backbone, low-capacity multitask head over
ESM-2 650M\cite{esm2}, with an explicit AlphaFold\cite{af2} structure channel
block and one linear read-out per task ($\sim$2M trainable parameters; Methods).
The pre-registered model of record ranks first of 115 entrants on three of five
CAID3 references and first among full-coverage entrants on all four CAID2
references, at coverage 1.00 throughout (Fig.~\ref{fig:model}b). The registered
primary family was fixed before the run: P1, beat AlphaFold-rsa --- confirmed
($+0.0097$, Holm-adjusted $p = 0.0252$); P2, beat PUNCH2 --- not significant
($+0.0043$, $p = 0.2022$). Superiority over PUNCH2 on pooled AUC is not
claimed.

On the calibration-invariant axis DisorderNet beats the CAID3 Disorder-PDB
winner on \emph{all five} benchmarks, per-target paired Wilcoxon, Holm-corrected
across the full family of thirty comparisons (Fig.~\ref{fig:model}a): Linker
$+0.1364$ ($p = 3.1\times10^{-5}$, 29 of 31 targets), Disorder-PDB $+0.0185$
($p = 1.4\times10^{-4}$), Binding $+0.1435$ ($p = 0.019$), Disorder-NOX
$+0.0496$ ($p = 0.041$), Binding-IDR $+0.2253$ ($p = 0.050$).

A temporal holdout tests generalisation without relying on any filter being
correct. Of 1{,}916 protein polymer entities released after the training caches
were built, 186 survive removal of exact training sequences and BLAST homologues
at $\ge 40\%$ identity --- 71\% of the ``new'' chains were already represented
(Fig.~\ref{fig:model}c), so a temporal cutoff alone is not a leak control. The
model of record reaches 0.8933 pooled and 0.8502 within-protein over all 186;
on the 61 chains where the training-free structural baselines can also be
computed it reaches 0.9095 and 0.8023 against AlphaFold-pLDDT's 0.8737 and
0.7816 (Fig.~\ref{fig:model}d, matched chain sets).

The holdout also settles a question CAID3 cannot. Five pre-registered
architectural variants were inseparable from their control on the benchmark, and
the capacity result says the benchmark could not have separated them whatever
they did --- so the natural objection is that real improvements are being
hidden. New data is not bounded by CAID3's capacity, and on these 186 unseen
chains \emph{none} of the four scorable variants beats its regime-matched
control, on pooled AUC ($-0.0033$ to $-0.0086$), pair-weighted within-protein
AUC ($-0.0047$ to $-0.0301$) or the unweighted form ($-0.0075$ to $-0.0137$).
There was nothing to resolve. And the unflattering result the set exists to
produce: the protein-bias variant is our best CAID3 model on Disorder-PDB and
the worst of three here, the only one failing to separate from a training-free
baseline. It buys CAID3 points and generalises worse.

% =====================================================================
\section*{Discussion}
% =====================================================================

The result we would most like to see used is the capacity bound, because it is
not about disorder. Any benchmark that ranks methods on annotations of finite
accuracy has a capacity, that capacity is $\lceil 1/2\eps \rceil$ regardless of
how many items it scores, and past it the ordering is not a function of the
observations. The converse is what separates this from a power calculation:
\code{unresolvable\_pair} constructs, for any two close methods, two ground
truths consistent with everything recorded, one favouring each. That is not a
limitation of a particular analysis; it is a statement that no analysis exists.

The obvious objection --- \emph{then nothing can be ranked} --- is answered by
the second result. Because a pair reverses only when both its residues flip in
opposite directions, the noise a pairwise protocol suffers is second-order where
a label protocol's is first-order. This is a general mechanism, not a fact about
proteins: it holds for any benchmark whose statistic is a function of ordered
pairs within groups. How much it buys must be measured, since the closed form
that would let a benchmark predict its own gain needs balanced classes.

The same shape appears once more, one step downstream. A screen's deflation
factor under arbitrary dependence is the harmonic number of the hypotheses it
states, that factor is attained and so cannot be argued away, and a region-level
report has the same false discovery proportion as its residue-level reading. The
unit of testing is a design choice exactly as the unit of scoring is, and
choosing it well is worth an order of magnitude in discoveries at no cost in
what is reported. Two independent instances of one principle: \textbf{where the
labels are noisy or the tests are many, what a study can resolve depends on a
unit that is usually chosen by convention rather than by argument.}

For CAID specifically we recommend reporting the within-protein component
alongside the pooled AUC, and printing a capacity --- the number of entrants the
reference can order, with everything below it an unresolved group rather than a
rank. CAID4's references are unpublished, so both are adoptable now.

\textbf{What we do not claim.} The closed form $\eps_{\text{pair}} \le 2\eps^2$
is proved but does not apply to CAID3, so the sixfold gain is measured and a
benchmark adopting the protocol must measure its own. The screen we ran is a
demonstration on a public benchmark, not a discovery about biology, and its
realised false discovery proportion is far below nominal, as arbitrary-dependence
control should be. DisorderNet is not state of the art on the metric CAID
reports: $+0.0043$ over PUNCH2 at $p = 0.20$, and no framing makes that a win.
The decomposition does not overturn the leaderboard ($\rho \ge +0.849$). The
pairwise ordering is reproducible but not shown to be \emph{more} valid. The
temporal and conformal analyses are exploratory; only the CAID3 primary families
were pre-registered. And the capacity result bounds the maximum size of a
certifiable family --- it is never promoted to a claim that such a family has
been certified.

Finally, several of this project's own results were wrong and were caught by
measurement rather than review --- including, while this manuscript was being
written, a headline capacity of 72 that formalising the denominator turned into
51, and a bound we had called tight whose hypothesis fails on the very data we
had checked it against. Each is documented with the measurement that caught it
(Supplementary Note~S5). A benchmark culture that reports only the number cannot
catch any of them.

% =====================================================================
\section*{Data availability}
% =====================================================================
CAID3 and CAID2 references and prediction archives are public at
\url{https://caid.idpcentral.org}; md5 checksums and composition counts of every
file used are in Supplementary Table~S12. MobiDB per-structure annotations are
public at \url{https://mobidb.org}. The temporal holdout is constructed from RCSB
\code{UNOBSERVED\_RESIDUE\_XYZ} annotations and is deterministic given the
cutoff date. All derived tables are provided as Supplementary Data
(Tables~S1--S15, comma-separated).

\section*{Code availability}
Model, evaluation pipeline, analysis scripts, the Lean~4 development and the
figure generator are available in the project repository. Pre-registrations are
committed with timestamps predating the runs they govern.

\section*{Author contributions}
T.R.M. conceived the study, developed the theory and the formal development,
implemented the model and the analyses, and wrote the manuscript.

\section*{Competing interests}
The author declares no competing interests.

\section*{Acknowledgements}
Computation was performed on the Advanced Research Computing at Hopkins
(ARCH) core facility.

% =====================================================================
\bibliographystyle{unsrt}
\bibliography{refs}

% =====================================================================
\clearpage
\section*{Figures}
% =====================================================================

\begin{figure}[htbp]
\centering
\includegraphics[width=\linewidth]{figures/fig1.pdf}
\caption{\textbf{CAID's statistic is a calibration contest.}
\textbf{a}, Share of the metric's comparison pairs that are within-protein
versus between-protein, for each CAID3 reference ($n = 233$, 178, 49, 42 and 31
targets). \textbf{b}, The rank of each benchmark's declared winner on the
within-protein axis, among full-coverage entrants. \textbf{c}, Inversion
certificates on Disorder-NOX: the measured between-protein gap divided by the
gap \code{inversion\_requires\_between\_gap} requires, log scale. Every
certificate holds.}
\label{fig:decomp}
\end{figure}

\begin{figure}[htbp]
\centering
\includegraphics[width=\linewidth]{figures/fig2.pdf}
\caption{\textbf{Capacity, and the way past it.}
\textbf{a}, Methods orderable, $\lceil 1/2\eps \rceil$, against annotation error
rate, with the measured $\eps$ marked and the 117-entrant line for scale.
\textbf{b}, Capacity under the residue-level and pairwise protocols, from rates
measured on 2{,}746 structure pairs. \textbf{c}, The hypotheses of
\code{nuPair\_le\_two\_eps\_sq}, checked on those pairs rather than assumed: the
balance condition holds on 54, both conditions on 32, and the bound holds on all
32 of those.}
\label{fig:capacity}
\end{figure}

\begin{figure}[htbp]
\centering
\includegraphics[width=\linewidth]{figures/fig3.pdf}
\caption{\textbf{CAID3 re-scored under the pairwise protocol.}
\textbf{a}, Pooled rank against pairwise rank, all 357 method--benchmark pairs,
with the identity line; $\rho = 0.844$. \textbf{b}, The twelve largest rank
changes; dot, pooled rank; arrowhead, pairwise rank. \textbf{c}, Entrants the
leader is separated from at the Holm-corrected margin, per benchmark; grey
denotes the eligible field. \textbf{d}, Cross-round reproducibility: Spearman
correlation between the CAID3 ranking and the held-out CAID2 ranking, over the
methods scorable in both rounds ($n = 24$, 24, 29 and 47 methods).}
\label{fig:rescored}
\end{figure}

\begin{figure}[htbp]
\centering
\includegraphics[width=\linewidth]{figures/fig4.pdf}
\caption{\textbf{DisorderNet.}
\textbf{a}, Per-target within-protein AUC margins of DisorderNet-Ensemble
against PUNCH2 and against AlphaFold-rsa; point, mean difference; bar, 95\%
bootstrap confidence interval over targets; $p$, two-sided Wilcoxon signed-rank
Holm-adjusted over all 30 comparisons; labels give wins/targets. \textbf{b},
Placement on both CAID rounds under each benchmark's own reported metric.
\textbf{c}, Construction of the temporal holdout. \textbf{d}, Performance on
structures released after the training caches were built, restricted to the 61
chains on which every method --- including the training-free structural
baselines --- can be scored, so the comparison is matched.}
\label{fig:model}
\end{figure}

\begin{figure}[htbp]
\centering
\includegraphics[width=\linewidth]{figures/fig5.pdf}
\caption{\textbf{The price of a guarantee.}
\textbf{a}, Realised miss rate against the fraction of protein that must be
flagged to achieve it, for all 70 full-coverage CAID3 Disorder-PDB methods at
$\alpha = 0.10$; median over 12 protein splits. \textbf{b}, The same guarantee
bought two ways --- one global threshold (filled) versus a per-protein quantile
(open) --- for the eight cheapest methods and AlphaFold-rsa; the number is the
difference, i.e.\ what protein-level calibration is worth operationally.
\textbf{c}, The comparison AUC cannot make.}
\label{fig:operating}
\end{figure}

\begin{figure}[htbp]
\centering
\includegraphics[width=\linewidth]{figures/fig6.pdf}
\caption{\textbf{The price of a screen, and where to pay it.}
\textbf{a}, The deflation factor required under arbitrary dependence against the
number of hypotheses stated, with Bonferroni's linear factor for comparison and
the residue-level and region-level counts for CAID3 Disorder-PDB marked.
\textbf{b}, The correction bought back by stating one hypothesis per region
rather than per residue, per reference, against the floor $\log b - 1$ that
\code{harmonic\_gain\_log} proves; labels give the ratio of testing thresholds.
\textbf{c}, Sharpness: on the constructed joint law, uncorrected
Benjamini--Hochberg realises $\mathbb{E}[\mathrm{FDP}] = qH_m$ exactly, so it
exceeds its nominal level from two candidates on.}
\label{fig:screen}
\end{figure}

\clearpage
% =====================================================================
%  Methods  --  Nature Methods places this after the references.
% =====================================================================
\section*{Methods}

\subsection*{Data and provenance}

All CAID3 numbers derive from the official reference FASTAs
(\url{https://caid.idpcentral.org/assets/sections/challenge/static/references/3/})
and the official prediction archive (\code{predictions.zip}, 76{,}480{,}806
bytes, 117 \code{.caid} files, aggregate md5
\code{cdef7ae6dc3dd6ec8055aacdb14a54a7}). Per-file checksums and per-target
composition counts are in Supplementary Table~S12. Composition is cross-checked
against CAID's dataset API, which serves those counts separately from the files
themselves; the evaluator refuses any reference whose composition does not
match. Of the 117 prediction files, 115 produce a score on any given reference,
which is the denominator of every rank reported here; the capacity tables use
117, since a method that entered is a method the benchmark was asked to place.

The dataset used is ``CAID3 v3'', the challenge as scored. The API also serves
an entry named plainly ``CAID3'' (185 proteins), a different and smaller set that
does not reproduce the published leaderboard. Three prediction files
(\code{FoldUnfold}, \code{NeProc-binding}, \code{NeProc-disorder}) leave the
score field blank on declined residues; these parse to NaN, are counted, and are
never silently dropped.

The check that licenses every other number is that each challenge's published
leader is recomputed through our own pipeline. All five reproduce to the
published precision: PUNCH2 0.9552 against 0.955, ESMDisPred-2PDB 0.8855 against
0.885, DisoFLAG-PB 0.7760 against 0.776, bindEmbed21IDR-rawGeneral 0.6407
against 0.641, IPA-AF2-Linker 0.8985 against 0.897. If a leader stops
reproducing, nothing computed afterwards is trusted until it is explained.

\subsection*{Annotation error rate}

$\eps$ is estimated two independent ways from MobiDB, neither of which uses
CAID's own references (which are not independent annotations of each other and
return exactly zero disagreements).

\emph{Context-dependent residues.} MobiDB's
\code{derived-missing\_residues\_context\_dependent-th\_90} annotation marks
residues missing in some deposited structures of a protein and observed in
others. Pooled over the 160 CAID3 Disorder-PDB targets with a usable record:
4{,}886 of 61{,}013 evidenced residues, $\eps = 0.0801$.

\emph{Structure pairs.} For each protein, every pair of its deposited structures
(\code{derived-missing\_residues-mobi-\{PDBID\}\_\{CHAIN\}}), restricted to
residues both structures cover, gives a direct disagreement rate. Over the 159
targets with $\ge 2$ structures and 2{,}746 structure pairs:
$\eps_{\text{label}} = 0.0651$ pooled and 0.0342 median protein. The same
enumeration gives $\eps_{\text{pair}} = 0.0100$ pooled and 0.0006 median, on the
denominator \code{card\_comparablePairs} defines --- the pairs both structures
order, $2(|T \cap L|\cdot|(T \cup L)^{c}| + du)$, and no larger set. The same run
records the balance hypothesis rather than assuming it: 32 of 2{,}746 structure
pairs satisfy both hypotheses of \code{nuPair\_le\_two\_eps\_sq}, and the bound
holds on all 32.

Per-structure coverage is approximated by the span of that structure's annotated
region, because MobiDB publishes per-structure \emph{missing} regions and not
per-structure coverage; residues outside a structure's span are treated as
undetermined rather than observed, which is the conservative direction, since
counting them as observed would manufacture disagreement from absence. Both
rates are worst-case budgets over structure pairs, not probability models,
matching the noise model the capacity theorem uses.

\subsection*{Formal development}

All theorems cited are Lean~4, \code{sorry}-free, and depend only on
\code{propext}, \code{Classical.choice} and \code{Quot.sound}. The development
covers:

\begin{itemize}\itemsep1pt \parskip0pt
\item the AUC decomposition and its invariance
(\code{auc\_pooled\_decomp};\ \code{auc\_within\_strictMono\_invariant};\
\code{auc\_pooled\_shift\_diff};\ \code{inversion\_requires\_between\_gap});
\item label noise and certification
(\code{LabelNoise.ranking\_certified};\ \code{RobustCertificate});
\item benchmark capacity in three score models
(\code{card\_le\_benchCapacity};\ \code{benchCapacity\_attained};\
\code{benchCapacity\_noise\_only};\ \code{unresolvable\_pair};\
\code{over\_capacity\_has\_close\_pair});
\item pairwise scoring
(\code{discordant\_eq\_flip\_product};\ \code{discordant\_le\_noise\_sq};\
\code{card\_comparablePairs};\ \code{noise\_orderKey};\
\code{pairwise\_capacity\_bound};\ \code{nuPair\_le\_two\_eps\_sq});
\item calibration
(\code{Calibration.risk\_decomposition};\ \code{validity\_is\_free});
\item multiple testing under dependence
(\code{wgt\_decomp};\ \code{selfConsistent\_fdr\_le\_harmonic};\
\code{bh\_fdr\_eq\_harmonic};\ \code{harmonic\_factor\_sharp};\
\code{fdp\_lift};\ \code{region\_screen\_fdr\_le};\
\code{harmonic\_gain\_log});
\item complexity (\code{biasThreshold\_hard};\
\code{crossed\_bound\_loose\_unbounded}).
\end{itemize}

A full inventory is Supplementary Note~S6.

Where the development supplies a worked example the analysis code reproduces it
exactly --- \code{Example.inversion} returns within 1 / pooled 2/3 against within
2/3 / pooled 5/6 --- which is the strongest available check that the Python
computes the statistic the theorems describe.

Two statements are used ahead of their formalisation and flagged as such:
\code{auc\_target\_strictMono\_invariant}, the per-target form of the AUC
invariance, which is what the protocol's step 3 needs (the published theorem is
stated for the pair-weighted mean, and the invariance holds target by target
before any averaging); and the validity of conformal risk control, cited rather
than formalised.

\subsection*{The pairwise protocol}

For a reference with targets $t$ and per-residue labels in $\{0,1,-\}$:

\begin{enumerate}\itemsep1pt \parskip0pt
\item \emph{Eligibility.} A method is scored iff it supplies a finite
prediction for every evaluated residue of every target, at the reference's
length. Declining targets raises a within-protein score, so coverage is a gate,
not a covariate. Ineligible methods are listed as \emph{not scored}, never as
last.
\item \emph{Per-target statistic.} For each target carrying both classes among
evaluated residues, the Mann--Whitney AUC over that target's residues alone.
Single-class targets contribute no ordered pair and are skipped, with the count
reported.
\item \emph{Method score.} The unweighted mean of the per-target AUCs. One
protein, one vote; pair-weighting lets a few long chains carry the number.
\item \emph{Ranking.} By that mean, descending.
\item \emph{Separation.} Two methods are reported as ordered iff a paired test
over targets rejects equality --- Wilcoxon signed-rank on the per-target
differences, Holm-corrected within the benchmark. Otherwise they are reported as
tied, at the same rank.
\item \emph{Capacity.} The number of methods the reference can order is
$\lceil 1/2\eps_{\text{pair}} \rceil$, with $\eps_{\text{pair}}$ estimated from
repeat determinations of the same protein. Methods beyond that are reported as an
unresolved group, never as a rank.
\end{enumerate}

The unweighted mean of step 3 differs from the pair-weighted
$\mathrm{AUC}_{\text{within}}$ that the decomposition identity requires. Both are
calibration-invariant, since the invariance holds target by target before any
averaging; the unweighted form is what a paired per-target design naturally uses
and is not dominated by long chains. A separation count is not a chain: the
capacity of step 6 bounds a certified totally ordered family, whereas $k$
separations against one leader are $k$ statements about the pairs
$\{\text{leader}, X\}$, which need not compose into a chain of $k{+}1$.

\subsection*{Model}

\emph{Backbone.} ESM-2 650M\cite{esm2}, frozen. A learned scalar mixture over
its 33 hidden layers (33 parameters) replaces fine-tuning. The design follows a
measurement: in the small-data regime here ($\sim$1M evidenced residues from
2{,}340 proteins) a 69.9M-parameter LoRA configuration reached pooled AUC 0.7454
while a gradient-boosted tree on hand-crafted features reached 0.7804 and
AlphaFold pLDDT alone reached 0.7906. The bottleneck was capacity/data fit, not
the backbone.

\emph{Head.} Four residual dilated-convolution blocks, hidden width 256,
dilations $(1,4,16,32)$ --- a 213-residue convolutional receptive field for the
same parameter count as the default $(1,2,4,8)$. The \emph{dependency span} is
global, not 213: the head's GroupNorm pools statistics over the whole length
axis, so every output position already depends on every input position. This is
verified directly (\code{lite\_head.measured\_dependency\_span}: on a
401-residue input, perturbing residue 0 moves the logit at residue 400).

\emph{Structure channels.} Relative solvent accessibility, pLDDT (scaled),
contact density, an availability flag, and the local gradient of solvent
accessibility, encoded into a small learned block. Missing structure is explicit
rather than imputed: an absent AlphaFold entry is not ``buried and confident'',
it is no information. Post-hoc fusion of the structural signal failed
($0.9581 \rightarrow 0.9554$); feeding structure as an input lets the trunk
learn when to trust it.

\emph{Read-outs.} One $1\times1$ convolution per task ($\sim$257 parameters
each) over the shared trunk, deliberately linear: anything deeper would let a
task with 15{,}683 positive residues (Linker) grow private capacity, which is
the failure mode the architecture exists to avoid. The shared trunk carries
disorder's 336{,}014 positive residues into the small tasks.

\emph{Windowing.} Training and inference use a 1{,}022-residue window at stride
511 with a raised-cosine taper on the overlap. Every window of a chain is
assigned to the same side of every split, keyed on the parent sequence.

\emph{Variants.} \code{mt\_windowed} (no protein bias, no private trunk; the
pre-registered model of record), \code{mt\_pbias} (adds a per-protein bias
term), and \code{DisorderNet-Ensemble} (equal-weight rank fusion of the
checkpoints sharing the fixed validation holdout, with membership decided on the
holdout and never on CAID3; both fixed in advance).

\subsection*{Leak control}

Splits are homology-clustered with BLASTp, clustered once over the union of
proteins so a protein appearing in two tasks cannot land in different folds for
each. Benchmark targets and their $\ge 40\%$-identity homologues are removed
from training. A fixed 5\% validation holdout is selected by a salted hash of the
parent sequence --- membership is a property of the sequence, not of the run ---
so that two runs with different leak filters remain comparable; the salt is a
module constant with no command-line override, and holdout membership is written
into every checkpoint. The homology filter refuses to run if BLAST is
unavailable rather than holding out nothing, because a validation number
computed with a paralogue in training is worse than none: it looks like a
measurement.

The temporal holdout is built from protein polymer entities first released on or
after 2026-08-11, the day after the later of the two training caches (DisProt
2026-08-08, MobiDB 2026-08-10). Labels are RCSB's own
\code{UNOBSERVED\_RESIDUE\_XYZ} annotation, not a re-derivation. X-ray and
cryo-EM only; NMR is excluded, since every SEQRES residue has coordinates in
every model.

\subsection*{Statistics}

Every p-value belongs to a family fixed in advance and Holm--Bonferroni\cite{holm}
is applied within it --- Holm rather than Bonferroni because it is uniformly more
powerful and assumes no independence, which these comparisons badly lack, sharing
both our predictions and their targets. The primary CAID3 family is fixed in the
pre-registration and encoded in the evaluator so it cannot be redefined after the
numbers land; everything outside it is secondary, corrected within its own
family, and labelled exploratory.

Paired comparisons resample \emph{proteins}, not residues, since residues within
a chain are heavily correlated and a residue-level interval is far too narrow.
All tests are two-sided; 10{,}000 resamples for confirmatory runs. Bootstrap
p-values use both tail boundaries and the Davison--Hinkley $+1$
correction\cite{davison}, so a perfect null returns 1.0 rather than 0 and $p$ is
never reported below $1/(B{+}1)$. Where a figure or table gives $p$, the test,
the family size and the correction are stated with it; exact $n$ (targets,
chains, structure pairs or methods) is given in every legend.

Coverage is a gate, not a covariate. Measured against a fully-covering yardstick,
declined targets are harder by $+0.035$ to $+0.125$ AUC and have median length
1{,}214--1{,}684 against 344 across all CAID3 targets (Supplementary Table~S9).
We predict every target on every benchmark and the evaluator aborts rather than
scoring a subset. Both rank fields are always reported --- rank among all scored
entrants (CAID's own accounting, which embeds the advantage) and rank among
full-coverage entrants (the equal-footing field). Reporting only the first
understates us; reporting only the second is cherry-picking. Neither alone is the
truth.

\subsection*{Conformal analysis}

Split conformal prediction\cite{vovk} and conformal risk control\cite{crc} as
implemented in \code{colab/conformal.py} (34 tests). Scores are rank-transformed
within each method, because conformal needs a consistent score and not a
probability. Calibration and evaluation are on disjoint halves of the targets,
median over 12 protein splits, with the interquartile range recorded. The
implementation refuses by name to report a guarantee at $\alpha < 1/(n{+}1)$,
which is arithmetic about $n$ rather than a fact about any method.

A per-residue promise is \emph{not} valid here, and the shortfall is measured
rather than hidden: residues within a chain are not exchangeable, so the split
is by chain, and at a 90\% target the realised per-residue coverage on the
temporal holdout is 0.876. It is quoted as a measurement of how far the
assumption is strained, not as a guarantee.

\subsection*{Screening analysis}

Two computations. \code{region\_screen.py} counts, for each CAID3 and CAID2
reference, the evaluated residues and the maximal same-label runs, and reports
both harmonic numbers, their difference, the proved floor $\log b - 1$, and the
ratio of the two testing thresholds. A run is broken by any unevaluated residue,
so a region interrupted by missing evidence counts as two candidates --- the
conservative direction, since it states more hypotheses rather than fewer.
Harmonic numbers are summed from the small end.

\code{conformal\_screen.py} runs the screen itself on CAID3 Disorder-PDB (304
usable chains, 98{,}938 evaluated residues). Candidates are consecutive blocks
of $b$ residues, the partition \code{stdBlocks} describes; a block's score is
the mean predicted disorder over its evaluated residues and a block is truly
disordered when the majority of those residues are, so that the truth and the
report are both unions of whole blocks, which is the condition \code{fdp\_lift}
needs. P-values are split conformal against the known-ordered blocks of a
calibration set of 80\% of chains,
$p = (1 + \#\{\text{calibration blocks with score} \ge s\})/(n_{\text{cal}}+1)$,
which is superuniform under exchangeability of the ordered blocks --- the only
assumption \code{selfConsistent\_fdr\_le\_harmonic} makes. It is also why the
correction must survive arbitrary dependence: every candidate is compared against
the \emph{same} calibration set, so the p-values are dependent by construction.
Four procedures are compared at the same guaranteed residue-level rate over 50
random splits (region-level BY at $\alpha/H_M$; residue-level BY at
$\alpha/H_n$; region-level uncorrected BH; region-level Bonferroni), and
\code{fdp\_lift} is checked rather than assumed by printing the region-level and
residue-level false discovery proportions of the same report side by side. The
screen is a demonstration on a public benchmark; no biological discovery is
claimed.

\subsection*{Reproduction}

Figures are regenerated from the analysis outputs by \code{paper/make\_figures.py}
and supplementary tables by \code{paper/make\_supplementary.py}; each panel reads
the JSON its job wrote, and the two panels that quote a table parse it from the
committed result document rather than restating it, so a figure cannot drift from
the number it illustrates. Job identifiers, batch scripts and raw outputs for
every analysis are listed in Supplementary Table~S13.

\subsection*{Reporting summary}

Further information on research design is available in the Nature Portfolio
Reporting Summary linked to this article.


\end{document}
